\section{The problem}

A corpus-based probabilistic grammar can predict held-out observations from its
own source corpus and still fail when it's moved to another corpus. External
validation needs more than a single discrimination statistic. It needs
probability scoring, calibration, comparison with target-native predictions on
the same observations, an account of which predictors carry the transport, and a
specification of the opportunity set over which the predictions are defined. This
paper assembles that regime and runs it on a canonical case. Throughout,
\term{probabilistic grammar} names the fitted production model, not a claim
about speakers' competence or a community's licensing of forms; external
validation bears on the model's predictive transport, not directly on either.

The English dative alternation is that case. Speakers can say \mention{give Kim
the book} or \mention{give the book to Kim}, and the choice is uneven: it varies
with verb, length, animacy, definiteness, pronominality, discourse status,
corpus, and register. The unevenness is easy to misread, because a rare pattern
can be lexically restricted, strongly dispreferred, or tied to a discourse
context without falling outside the grammar. \textcite{bresnan2007} modelled that
conditioning and argued that corpus probabilities can reveal grammatical
possibilities informal judgements miss.

The model has since been shown to travel: trained on one variety and scored on
another, it keeps much of its accuracy
\citep{BresnanHay2008,KendallBresnanVanHerk2011,EngelEtAl2025}. The open question
isn't whether it travels but how well, once transport is treated as a full
predictive-evaluation problem, and what a fitted production probability can and
can't say about grammar. That last is a question of \term{projectibility}, what
observing a profile in one setting licenses about another \citep{Goodman1955}.

The argument proceeds by separating validation claims that are often collapsed.
Reproducing the source model asks whether the classic production profile is
stable in its original data. Scoring frozen predictions in BNC2014 asks whether
that profile carries predictive information out of domain. Comparing those
predictions with a target-native model on the same rows asks how much local
structure the source model misses. Calibration and ablation then ask whether the
transport is a matter of ranking, probability level, lexical base rates, or
non-lexical conditioning. Finally, an acceptability bridge asks where a
production probability is licensed, and where it isn't, when carried into a
judgement channel.
